\section{绪论}

\subsection{具身智能体在动态自然环境中的思考与行动权衡}

携带先验环境底图的具身移动智能体在未知或动态演化的自然环境中执行探索任务时，必须在“深入思考”与“立即行动”之间持续做出资源分配决策\cite{horvitz1987reasoning, russell1991right, shen2025thinking}。与运行在充裕算力服务器上的纯软件智能体不同，部署于野外的自主移动平台受制于严格的机载电池能量上限、受限的嵌入式计算速率以及不可逆的物理作业时间窗口\cite{xiao2020thinking, magnino2026metareasoning}。当智能体依托大规模历史观测构建的先验底图进入真实物理场时，先验底图虽能提供大尺度空间结构线索，却不可避免地因环境演化与观测盲区而存在时效性滞后与小尺度失真\cite{jiang2015horizon}。智能体若仅凭粗糙的瞬时规划仓促行动，极易在复杂的局部梯度中迷失方向，错失核心目标区域；反之，若频繁唤醒高功耗加速计算模块在内部底图模型上开展高强度轨迹推演，或关停推进系统在原地反复采样以压低局部感知噪声，则会直接挤占用于长距离空间机动的宝贵时间与推进电能\cite{cuadron2025overthinking, elboher2023formal}。因此，如何在算力受限、能量受限与底图失真的三重耦合约束下，科学界定内部规划推演与外部物理采样的收益边界，构成了实时科学计算与智能导航交叉领域的核心科学问题\cite{hunt2021formalizing}。

在并发控制架构下，智能体的机体机动与大脑规划虽可在同一时间步内并行执行，但两者的资源竞争关系依然通过电能消耗与决策延迟紧密交织\cite{xiao2020thinking, elboher2023formal}。一方面，维持基础巡航仅需后台低功耗处理器以较低频率更新路径，而开启嵌入式GPU进行高吞吐量轨迹搜索则会显著抬升整机功耗，直接缩短电池耗尽前的总航行里程；另一方面，若智能体选择关停推进电机以降低航行功耗并提高局部水体采样密度，虽能节约推进能量，却牺牲了在有限时间窗口内横跨大尺度空间寻找全局极值的机动能力。更为深刻的是，智能体的“思考”完全建立在其内部持有的环境概率模型之上：当内部模型存在固有的先验不确定性时，单纯增加规划器的搜索算力能否真正转化为物理世界中的采样收益，取决于规划器的搜索深度是否已触及先验底图的信息上限\cite{jiang2015horizon, smith2006curse}。

\subsection{海洋大数据驱动的东海赤潮应急巡测场景}

将上述通用决策问题置于海洋科学应用场景，卫星遥感大数据支撑下的东海有害赤潮应急机动巡测提供了一个极具代表性的真实试验场\cite{fossum2019, lou2014diurnal}。有害赤潮是海洋浮游植物在特定水文气象条件下暴发性增殖引起的水体变色与生态灾害现象；根据自然资源部发布的海洋灾害公报，我国海域仅2023年即发现赤潮46次，累计面积约1500平方千米，其中有毒有害赤潮达29次\cite{mnr2024bulletin}。在浙江近海与长江口外海域，晚春至初夏时节频繁暴发以东海原甲藻为优势种的大规模赤潮，高密度藻华斑块的叶绿素a浓度（反映水体浮游植物生物量与赤潮密度的核心光学指标）常超过$6\text{ mg m}^{-3}$，并在潮汐与表层海流驱动下发生快速的日际漂移与形态演变\cite{lou2014diurnal}。在赤潮应急响应中，岸基指挥中心亟需派遣自主水下航行器（AUV，一种无需母船实时缆控、依靠机载能源与智能算法自主航行采样的水下机器人）深入暴发海域，在24小时黄金应急窗口内尽早定位叶绿素浓度最高的核心水体，并随时响应岸基中心的位置查询\cite{fossum2019, remus100}。

\begin{figure}[htbp]
  \centering
  \includegraphics[width=\linewidth]{fig_overview}
  \caption{海洋大数据驱动的AUV赤潮应急巡测与“思考—行动”资源分配框架总览。中央主视图展示东海测试任务窗口（2025年4月4日）的三维海洋叶绿素场、驶向赤潮核心区（黄色圆环）的AUV及其轨迹（蓝色为低功耗后台规划巡航段，橙色为唤醒GPU深度思考段），海面上方悬浮6个假想海洋切片（4个深蓝边框样本供规划器搜索，2个橙色边框样本留出用于无偏评估）及提供先验底图的遥感卫星。侧边子图：（A）八年卫星叶绿素遥感序列及底图随滞后天数的退化过程；（B）航行器内部48个假想海洋的搜索集（32个）与留出集（16个）解耦架构，以及出发前对五种时间分配策略的闭环预演机制；（C）每20分钟决策周期内的机体与算力选项、电池功耗分配及向岸基指挥中心实时输出的最高浓度水体位置估计。}
  \label{fig:overview}
\end{figure}

卫星海洋水色遥感大数据为AUV巡测提供了不可或缺的大尺度先验视野，但其固有的时空分辨率限制与时效性滞后决定了现场自主观测的不可替代性\cite{beckers2003dineof, ye2021cocts}。由于东海春夏季节云雨频繁，原始光学卫星影像存在大面积云层遮挡，业务化海洋大数据产品通常采用经验正交函数缺失数据插值重构算法（一种利用时空矩阵低秩特征填补云层盲区的统计重构方法）生成空间完整、分辨率为2千米的逐日融合叶绿素底图\cite{beckers2003dineof}。然而，如图\ref{fig:overview}所示，当AUV抵达现场时，受卫星过境周期、地面反演处理及插值平滑滞后的影响，机载装订的往往是1至3天前的滞后底图；加之重构算法抹平了次中尺度锋面与小尺度丝状藻华结构，且近岸浑浊水体反演存在显著本底偏差\cite{ye2021cocts}，真实海洋的赤潮核心位置与滞后卫星底图之间往往存在十余千米的空间偏移与显著的幅值涨落。因此，AUV必须将卫星大数据转化为概率化的先验认知，一边通过机载荧光计实时测量沿途水体叶绿素浓度以修正内部底图，一边动态权衡算力、电能与航时分配。

\subsection{现有元推理与自适应采样方法的局限及本文贡献}

面对机载算力与采样行动的资源分配难题，直接套用直观的启发式规则、经典单步元推理或离线强化学习策略均会因海洋环境与规划器的内在特性而失效。第一种直观做法是预设单一的固定分配规则（例如全程低功耗巡航，或按固定比例唤醒GPU思考），但固定规则无法兼顾不同硬件平台的算力差异：在算力孱弱的平台上全程巡航会因轨迹搜索不充分而陷入局部次优，而在算力充裕的平台上频繁开启高功耗GPU则会白白消耗推进电能并缩短有效航时。第二种做法是采用人工智能领域的经典单步元推理框架\cite{horvitz1987reasoning, russell1991right, lin2015metareasoning, sung2021quit}，即在每个决策步在线比较“多思考一步的计算价值”与“多采样一步的信息价值”；然而，经典方法直接利用规划器在搜索集合上的内部得分跃升来估计计算价值，忽略了连续轨迹优化在高预算下对有限样本集合的严重过拟合——即决策分析中的“优化器诅咒”\cite{smith2006curse, vanhasselt2010double}，导致规划器在真实收益早已饱和后仍误以为继续思考大有裨益，从而陷入持续消耗电能的“过度思考”陷阱\cite{cuadron2025overthinking}。第三种做法是从历史巡测任务结果中直接训练策略选择网络\cite{magnino2026metareasoning}；但在真实赤潮场景中，不同任务窗口间的自然环境复杂度差异（遗憾值标准差约0.08）远大于不同优良策略之间的平均性能差距（通常不超过0.01），极低的信噪比使得离线学习极易过拟合历史环境噪声。此外，试图单纯依靠深度学习时空预测模型外推卫星底图以消除先验滞后的做法\cite{assran2025vjepa2}，同样因日际小尺度湍流与生物生长过程的随机性而无法超越持续性基线。

为突破上述瓶颈，本文跳出“依赖内部搜索轨迹做单步短视评估”与“依赖历史任务结果做离线拟合”的传统定式，提出利用海洋大数据校准不确定性、通过搜索与评估集合解耦消除过思考偏差，并在出发前开展全任务周期假想预演的自校准元控制方法。本文的核心贡献概括如下：
\begin{enumerate}
  \item \textbf{揭示了连续轨迹规划中的过思考偏差机理并提出搜索—留出双集合解耦架构：} 将AUV内部维护的48个假想海洋严格划分为32个搜索样本与16个留出验证样本；实证发现当规划推演预算跨越四个数量级增长时，规划器在搜索集上的自评得分持续虚高，而留出集评估值与真实海洋收益之差始终稳定在0.005以内，从根本上解决了有限集合规划中的价值高估难题。
  \item \textbf{阐明了海洋移动采样中“思考—巡航—停航重采样”的硬件与环境相图规律：} 基于东海八年3197天遥感大数据与422个赤潮任务窗口，定量刻画了卫星底图时效性退化图谱，证明了深度学习时空外推无法替代现场观测，并揭示出高功耗GPU深度思考仅在基础算力不足以覆盖有效规划视野时产生净收益，而巡航途中停航漂移重采样在中尺度相关场中始终得不偿失。
  \item \textbf{提出了基于留出假想海洋的行前闭环预演元控制器：} 在任务出发前以留出假想海洋为合成真实环境，结合真实机载算力速率与电池功耗模型，采用共同随机数对候选时间分配策略开展24小时全周期闭环预演；该方法无需任何针对特定硬件的历史任务调参，即可在不同算力平台上自动逼近事后最优调参策略，并显著优于单一固定规则与经典单步元推理。
\end{enumerate}

\section{东海赤潮遥感大数据特征与先验不确定性量化}

\subsection{八年多源融合叶绿素遥感数据集与赤潮任务窗口构建}

构建高保真度的海洋自主采样仿真基准，必须依托长时序、无缝隙且具备真实地球物理统计特征的卫星遥感大数据。由于单颗极轨水色卫星受刈幅宽度与云层遮挡限制难以提供连续的每日海洋表面视图，本文选用美国国家海洋和大气管理局海岸观测网发布的每日多源融合叶绿素a浓度遥感产品作为核心海洋大数据底座（公共领域开放许可，数据集标识为\texttt{noaacwNPPN20S3ASCIDINEOF2kmDaily}）。该数据集融合了搭载于苏奥米国家极地轨道伙伴卫星与美国国家海洋和大气管理局20号卫星上的可见光红外成像辐射仪，以及搭载于哨兵三号甲星上的海洋陆地颜色仪观测数据，并采用经验正交函数缺失数据插值重构算法对云层遮挡区域进行时空一致性重构\cite{beckers2003dineof}。数据空间分辨率为2千米，时间分辨率为逐日，时间跨度覆盖2018年1月1日至2026年10月23日共3197天；研究区域截取东海核心海域（北纬$25^\circ\text{--}33^\circ$，东经$119.5^\circ\text{--}127^\circ$），每日构成$385\times 361$个空间网格像素。同时，为检验多模态海洋环境变量对赤潮演变的辅助预测能力，本文同步引入美国喷气推进实验室发布的多尺度超分辨率海表温度每日无云产品（版本4.1，数据集标识为\texttt{jplMURSST41}）\cite{chin2017mur}，将其从原始$0.01^\circ$分辨率重采样至$0.04^\circ$以匹配水色网格。因海洋叶绿素a浓度在空间上呈典型的对数正态分布，跨越数个数量级，本文所有叶绿素数值、误差统计与遗憾值指标均在以10为底的对数空间（即$\log_{10}(\text{叶绿素a浓度}/\text{mg m}^{-3})$）下定义与计算。

为确保实验评估聚焦于具有真实应急响应价值的离岸赤潮暴发事件，本文制定了严格的地球物理筛选准则，从八年春季至初夏序列中提取标准化的任务窗口。具体而言，首先在每年4月至6月的赤潮高发期内，搜索浙江近海陆架区域（北纬$26.5^\circ\text{--}31.5^\circ$，东经$121.5^\circ\text{--}125^\circ$）离岸距离至少30千米处、经空间平滑后局部极值超过$5\text{ mg m}^{-3}$的显著藻华中心；随后以该极值点为基准，施加最高达16千米的随机空间偏移以防止目标恒定处于区域几何中心，进而截取$80\text{ km}\times 80\text{ km}$（对应$40\times 40$个2千米网格单元）的正方形海域作为单次巡测任务窗口。为保障AUV航行安全并确保场内具备有意义的浓度梯度，将距离陆地或岛屿岸线15千米以内的网格标记为禁航区，仅保留可航行水域占比不低于80\%且窗内叶绿素对数浓度极差（对比度）不低于0.3的任务窗口。按此标准共筛选出422个高质量赤潮任务窗口；为严格杜绝跨年度数据泄露，按时间顺序将2018至2022年的274个窗口划分为开发集（用于先验误差统计校准与基线调参），将2023至2026年的148个窗口划分为完全独立的测试集。在每次任务开始时，AUV随机部署于任务窗口边界上的任一可航行网格点。

\subsection{卫星底图时效性退化图谱与小尺度海真偏差建模}

卫星遥感底图在AUV下水时刻的先验不确定性主要源于两个相互独立的物理过程：一是底图获取至任务执行期间（通常滞后1至3天）水体平流与生物生消引起的时效性退化，二是遥感反演与插值平滑相对于同日真实表层海水的次网格尺度表征偏差。通过对全部3197天东海叶绿素遥感序列进行滑动差分统计，本文绘制了3天滞后叶绿素对数浓度变化量的季节性空间分布图谱（图\ref{fig:atlas}）。统计结果显示，东海叶绿素场的3天变化剧烈程度具有显著的季节节律与空间非均质性：全海域中位数标准差在春季最高，达到0.127；夏季次之，中位数为0.107，但其空间分异极为强烈，全海域90\%分位数高达0.204，在浙闽沿岸流锋面与长江冲淡水羽流外缘形成了一条标准差高达0.25至0.30的强动态变化带；秋季与冬季则逐步减弱，全海域中位数分别降至0.103与0.081。在274个春季开发集赤潮任务窗口内部，真实目标日底图减去$k$天前滞后底图的时效性退化误差统计表明，当滞后天数$k=1, 2, 3$天时，误差标准差分别为$s_1=0.125$、$s_2=0.176$与$s_3=0.210$，空间自相关特征尺度约为$\ell=5$个网格单元（即10千米），且伴随$+0.016$、$+0.032$与$+0.045$的微弱正均值偏置（反映了赤潮暴发初期藻华生物量整体处于净增长阶段的统计特征）。

\begin{figure}[htbp]
  \centering
  \includegraphics[width=\linewidth]{fig_atlas}
  \caption{东海海域（北纬$25^\circ\text{--}33^\circ$，东经$119.5^\circ\text{--}127^\circ$）八年3197天卫星叶绿素对数浓度3天变化幅度的四季气候态分布图谱。陆地区域以灰色标示；春季子图中叠加了全部422个赤潮巡测任务窗口的中心位置，其中蓝色圆点为2018至2022年开发集窗口（274个），绿色圆点为2023至2026年测试集窗口（148个）。浙闽沿岸与长江冲淡水锋面在春夏两季呈现出显著的高动态变化特征。}
  \label{fig:atlas}
\end{figure}

除时间滞后外，即便是目标日当天的融合卫星底图也并非真实海洋的完美写照。一方面，经验正交函数插值重构算法在填补云区时截断了高阶奇异值，使重构场在时空上呈现明显的低通平滑特性\cite{beckers2003dineof}；另一方面，针对海洋一号丙星水色水温扫描仪在东海与南海叶绿素反演产品的现场比对验证表明，卫星遥感反演值与现场实测值之间存在约39\%的无偏相对百分比差异（折算至以10为底的对数空间约为0.15的均方根偏差）\cite{ye2021cocts}。为在仿真环境中忠实复现这一遥感底图无法分辨的次中尺度斑块结构，本文将任务窗口目标日的真实海洋叶绿素场$c(\mathbf{x})$建模为目标日卫星底图叠加一个空间相关的零均值高斯随机场，其基准标准差设为$s_{\text{sat}}=0.15$，空间相关尺度设为3个网格单元（6千米）；同时在机理相图实验中，将海真相对于卫星底图的偏离程度在$0.05$至$0.30$范围内开展全参数扫描，以全面检验不同遥感精度下的决策规律。

\subsection{深度时空世界模型外推的局限性与现场观测的不可替代性}

面对滞后1至3天的卫星底图，人工智能赋能科学计算领域最直观的改进思路是利用多年历史遥感序列训练一个深度时空预测模型（即“海洋世界模型”），在AUV下水前直接将滞后底图向前推演至任务当日，从而在源头上消除底图的时效性退化。为严格检验这一技术路线在真实东海赤潮场景中的有效性，本文构建了一个参数量为550万的联合嵌入预测架构网络\cite{assran2025vjepa2}，以连续多日叶绿素a浓度场与同步海表温度场\cite{chin2017mur}作为多模态输入，利用2018至2021年每年3月至6月的历史序列进行训练，并在2022年验证集上进行超参数优选与早停控制，旨在预测未来1至3天的叶绿素场演变。

然而，在2023至2026年独立测试集上的系统评估给出了明确的否定性结论：深度时空世界模型对未来1至3天重构叶绿素场的预测精度未能超越最简单的“持续性基线”（即直接将最新可用的滞后底图作为未来预测值）。具体而言，在提前1天、2天与3天的预测任务中，深度学习模型的测试集均方根误差分别为0.1229、0.1439与0.1487，而直接沿用滞后底图的持续性基线均方根误差分别为0.1226、0.1439与0.1474；即使在2022年验证集上，深度学习模型相对于持续性基线的误差降幅也仅为微不足道的1.4\%，且其预测集合的离散度与真实误差之比仅为0.70至0.93，存在明显的方差低估倾向。这一结果具有深刻的物理与信息论根源：首先，业务化插值重构底图本身已是经时空低秩滤波后的平滑场，其中可由大尺度历史轨迹确定性外推的低频趋势分量在1至3天尺度上变化极慢，持续性假设已近似最优；其次，真正决定赤潮峰值斑块日际位移的10千米级次中尺度涡丝、潮致锋面摆动以及藻华非线性暴发增长，受控于未观测的局地风应力与营养盐脉冲，在粗分辨率日均遥感序列中属于不可预测的随机创新过程。因此，试图依靠离线神经网络外推来恢复小尺度赤潮峰值位置在物理上是不可行的；海洋大数据的正确打开方式，是从历史序列中精确校准先验误差的幅值与空间相关结构，进而指导AUV携带传感器深入现场开展自适应同化观测。

\section{随流参考系下的并发采样与能耗受限决策问题建模}

\subsection{随流漂移参考系与二十四小时随时遗憾值评价指标}

在潮汐与表层海流强劲的东海陆架区规划AUV巡测路径，若直接在固定于海底的欧拉地理参考系下建模，不仅需要高精度的实时三维海流预报场，还会使规划器耗费大量算力去处理水体整体的刚性平流运动。地球静止海洋水色成像仪对东海赤潮日内演变的高频追踪研究表明，表层浮游植物藻华斑块在日尺度上主要作为被动示踪物随表层水体同步漂移\cite{lou2014diurnal}。基于这一物理特性，本文将整个巡测运动学、环境场表征与数据同化过程完全建立在随表层水体同步漂移的拉格朗日参考系（下称随流漂移参考系）中。在该参考系下，赤潮叶绿素场相对于水体保持准静止状态，AUV以恒定的对水航速机动，无需为“追逐”随流漂移的藻华而与海流做无谓对抗；海流平流与剪切变形对先验认知的影响，则自然地内化于卫星底图从过境时刻至任务时刻的时效性退化误差统计之中。

在$40\times 40$个网格单元（分辨率2千米，总面积$80\text{ km}\times 80\text{ km}$）构成的任务窗口$\mathcal{X}$内，设真实海洋在网格点$\mathbf{x}\in\mathcal{X}$处的叶绿素a对数浓度为$c(\mathbf{x})$，全局真实最高浓度为$\max_{\mathbf{x}\in\mathcal{X}} c(\mathbf{x})$。在总时长为$T=24\text{ h}$的应急响应窗口内，岸基指挥中心可在任意时刻$t\in[0, T]$向AUV查询当前已探明的最密集赤潮水体位置。为防止AUV仅凭未经证实的先验猜测上报虚假极值，任务规程要求AUV必须从截至时刻$t$已实际造访并采样的网格集合$\mathcal{V}(t)\subseteq\mathcal{X}$中，选取当前内部同化后验均值$\bar{c}_t(\mathbf{x})$最高的网格点作为实时回报位置$\hat{\mathbf{x}}(t)$：
\begin{equation}
  \hat{\mathbf{x}}(t) = \arg\max_{\mathbf{x}\in\mathcal{V}(t)} \bar{c}_t(\mathbf{x})
  \label{eq:answer}
\end{equation}
若AUV在$24\text{ h}$结束前因电池能量耗尽而终止机动，其已访集合$\mathcal{V}(t)$与回报位置$\hat{\mathbf{x}}(t)$在剩余时间内保持冻结。为综合衡量AUV“尽早发现高浓度核心区”与“最终定位精度”的双重应急效能，本文定义24小时时间平均随时遗憾值（下称随时遗憾值）$R$作为核心评价指标：
\begin{equation}
  R = \frac{1}{T}\int_0^T \left[\max_{\mathbf{x}\in\mathcal{X}} c(\mathbf{x}) - c\big(\hat{\mathbf{x}}(t)\big)\right] dt
  \label{eq:regret}
\end{equation}
该指标度量了在整个24小时应急窗口内，真实全局最大叶绿素对数浓度与AUV实时回报位置处真实叶绿素对数浓度之差的时间积分均值（越低越好）。此外，本文同步记录任务结束时刻（$t=24\text{ h}$）的终点遗憾值，以及终点回报位置真实浓度落入全局极值$0.1$个对数单位以内（在线性浓度尺度上相当于偏差不超过26\%）的命中率。

\subsection{航行器机电功耗模型与二十分钟并发决策周期}

为真实刻画移动观测中的能量与算力耦合约束，本文依照百米级轻型AUV平台（以REMUS 100型平台技术参数为基准）\cite{remus100}与嵌入式边缘计算模块（以Jetson Orin NX级GPU模块为基准）\cite{jetsonorinnx}构建机电功耗模型。航行器搭载总容量为$E_{\max}=1\text{ kWh}$（即$1000\text{ Wh}$）的锂离子电池组，额定对水巡航速度为$v=1.5\text{ m/s}$（约3节）。在巡航状态下，整机基础功耗为$45\text{ W}$（对应约22.2小时的纯巡航续航时间），由两部分构成：一是无论航行或漂移均需持续供电的$P_{\text{hotel}}=15\text{ W}$基础航电与传感器功耗（含导航、声学通信、荧光计及低功耗后台处理器），二是驱动螺旋桨以$1.5\text{ m/s}$航行的$P_{\text{prop}}=30\text{ W}$推进功耗。在计算单元方面，低功耗后台处理器常驻运行，其功耗已包含在$15\text{ W}$基础航电中，每分钟可执行$r_{\text{base}}\in[5, 200]$次轨迹推演（取决于具体硬件配置）；当需要深度规划时，航行器可唤醒嵌入式GPU加速模块，其推演吞吐量提升至后台处理器的50倍（即$50\,r_{\text{base}}$），但在基准配置下需额外消耗$P_{\text{gpu}}=20\text{ W}$电能（在高效能专用加速器扩展配置下额外功耗为$5\text{ W}$）。

整个24小时任务被划分为72个时长$\Delta t = 20\text{ min}$（即$1/3\text{ h}$）的离散并发决策周期\cite{xiao2020thinking, elboher2023formal}。在每个20分钟周期开始时，元控制器需同时为“机体”与“大脑”下达动作指令：
\begin{itemize}
  \item \textbf{机体动作选项：}（1）\textbf{巡航采样}，即开启推进电机（$30\text{ W}$），沿当前生效的规划路径以$1.5\text{ m/s}$航行$18\text{ km}$（跨越约9个2千米网格），机载荧光计每航行$1\text{ km}$采集1个独立样本（单周期共采集18个样本，每个2千米网格内平均获得2至3次独立测量）；（2）\textbf{停航漂移重采样}，即关停推进电机（推进功耗降至$0\text{ W}$，仅保留$15\text{ W}$基础航电），在随流漂移参考系下驻留于当前网格位置，荧光计每$4\text{ min}$采集1个独立样本（单周期共采集5个原位样本），以通过多次重复测量压低局部传感器随机噪声。每个独立荧光计样本的测量方程为：
  \begin{equation}
    y = c(\mathbf{x}) + \varepsilon, \quad \varepsilon \sim \mathcal{N}(0, \sigma_{\text{obs}}^2)
    \label{eq:obs}
  \end{equation}
  其中测量噪声标准差为$\sigma_{\text{obs}} = 0.1$。
  \item \textbf{大脑算力选项：}（1）\textbf{后台常规划}，即仅依靠低功耗处理器在20分钟内执行$20\,r_{\text{base}}$次轨迹推演（额外功耗$0\text{ W}$）；（2）\textbf{深度思考规划}，即唤醒GPU在20分钟内执行$1000\,r_{\text{base}}$次轨迹推演（额外功耗$P_{\text{gpu}}$）。
\end{itemize}
计算与机动在周期内完全并发运行：规划器在当前20分钟周期内读取最新同化后的假想海洋集合，并以航行器在当前周期结束时刻的预计到达位置作为新轨迹的起点（即精确时延补偿），搜寻最优后续路线；待当前20分钟周期结束时，新规划出的路线无缝替换旧路线并投入下一周期执行。当累计消耗电能达到$1\text{ kWh}$或任务时间满24小时时，任务立即终止。

\section{基于留出假想海洋的集合同化与行前预演元控制方法}

\subsection{数据校准的假想海洋先验集合与序列局地化同化更新}

为使AUV在缺乏确定性未来底图的条件下具备概率化的空间推理能力，本文不再让航行器维护单一的静态点估计底图，而是围绕滞后$k$天的卫星底图$c_{\text{map}}(\mathbf{x})$构建由$N=48$个“假想海洋”（即先验高斯随机场实现样本）组成的集合记忆库。基于第2.2节从八年遥感大数据中校准得到的误差统计量，第$i$个假想海洋样本$c^{(i)}(\mathbf{x})$在任务初始时刻（$t=0$）按如下公式生成：
\begin{equation}
  c^{(i)}(\mathbf{x}) = c_{\text{map}}(\mathbf{x}) + \sigma_k\, z^{(i)}(\mathbf{x}), \quad i = 1, 2, \dots, 48
  \label{eq:prior}
\end{equation}
式中，$\sigma_k = \sqrt{s_k^2 + s_{\text{sat}}^2}$为融合了$k$天底图滞后误差标准差$s_k$与海真小尺度偏离标准差$s_{\text{sat}}=0.15$的合成先验不确定性幅值（例如当底图滞后$k=3$天时，$s_3=0.210$，合成标准差$\sigma_3 \approx 0.258$）；$z^{(i)}(\mathbf{x})$为具有零均值、单位方差及平方指数空间协方差结构的独立同分布高斯随机场，其空间相关长度设为校准值$\ell = 5$个网格单元（10千米）：
\begin{equation}
  \text{Cov}\big(z^{(i)}(\mathbf{x}), z^{(i)}(\mathbf{x}')\big) = \exp\left(-\frac{\|\mathbf{x} - \mathbf{x}'\|_2^2}{2\ell^2}\right)
  \label{eq:cov}
\end{equation}
这48个假想海洋共同编码了航行器对当前海域赤潮空间分布的先验信念：每个样本本身都是一幅具有逼真次中尺度斑块结构的完整可能海洋，而48个样本之间的离散度则精确反映了先验底图的不确定性空间分布。

当AUV在巡航或漂移过程中获取到现场荧光计观测序列时，必须实时将点位测量信息传播至周围未观测网格并压缩集合不确定性。若采用标准蒙特卡洛集合卡尔曼滤波直接更新包含1600个网格点的状态向量，一方面向观测值人为注入随机扰动会引入额外的采样方差，另一方面仅用48个样本估计1600维协方差矩阵会产生严重的远距离伪相关，导致几十千米外的无关网格发生非物理跳变。因此，本文采用确定性的序列集合调整卡尔曼滤波算法\cite{anderson2001eakf}，并耦合五阶分段有理局地化函数\cite{gaspari1999correlation}对每个独立观测$y$进行逐一序列同化。具体而言，当在位置$\mathbf{x}_{\text{obs}}$获得观测$y$时，首先计算集合在该网格处的先验均值$\bar{c}_{\text{obs}}$与先验方差$\sigma_{\text{prior}}^2$，利用卡尔曼方差缩减公式精确调整该点处的集合离散度以匹配理论后验方差$\sigma_{\text{post}}^2 = (\sigma_{\text{prior}}^{-2} + \sigma_{\text{obs}}^{-2})^{-1}$，得出观测点处第$i$个样本的状态增量$\Delta c_{\text{obs}}^{(i)}$；随后，通过集合样本回归系数及以$2\ell$（10个网格单元，即20千米）为半宽尺度的局地化截断权重$\rho(\|\mathbf{x}-\mathbf{x}_{\text{obs}}\|_2 / 2\ell)$，将观测增量平滑映射至邻近网格：
\begin{equation}
  c_{\text{post}}^{(i)}(\mathbf{x}) = c_{\text{prior}}^{(i)}(\mathbf{x}) + \rho\left(\frac{\|\mathbf{x}-\mathbf{x}_{\text{obs}}\|_2}{2\ell}\right) \frac{\text{Cov}_{\text{ens}}\big(c(\mathbf{x}), c(\mathbf{x}_{\text{obs}})\big)}{\sigma_{\text{prior}}^2} \Delta c_{\text{obs}}^{(i)}, \quad i = 1, \dots, 48
  \label{eq:eakf}
\end{equation}
上述同化机制无需对高维矩阵求逆，单次观测更新耗时仅亚毫秒级，既保证了48个假想海洋在每次观测后严格满足贝叶斯后验均值与方差一致性，又将观测信息严格限制在物理合理的空间相关半径之内。

\subsection{轨迹规划中的优化器诅咒与搜索—留出双集合解耦机制}

在每个20分钟决策周期内，规划器的任务是从下一周期起点出发，在连续二维空间中搜寻一条能够在剩余任务时间$T_{\text{left}}$内最大化高浓度水体发现收益的后续航行路线。为在保持轨迹几何灵活性的同时控制搜索空间维度，本文将候选路线参数化为由当前预计终点$\mathbf{x}_0$依次连接两个中间航路点$(\mathbf{w}_1, \mathbf{w}_2)\in\mathcal{X}^2$构成的连续折线轨迹（共4个连续坐标自由度）。将一条候选路线在某一个假想海洋样本$c^{(i)}$上完整推演一次称为一次“推演”（rollout）：假设航行器沿该折线以巡航速度$v$航行，记推演时刻$\tau\in[0, T_{\text{left}}]$沿途已遇见的最高叶绿素浓度为$m_i(\tau)$（其初值$m_i(0)$设为当前已访网格集合上的最高后验均值估计，且一旦剩余可用巡航电池时间耗尽，$m_i(\tau)$即保持冻结），则该候选路线在第$i$个假想海洋上的路线价值$V_i$定义为剩余窗口内的时间积分均值：
\begin{equation}
  V_i(\mathbf{w}_1, \mathbf{w}_2) = \frac{1}{T_{\text{left}}}\int_0^{T_{\text{left}}} m_i(\tau)\, d\tau
  \label{eq:route_value}
\end{equation}
式\eqref{eq:route_value}将全局随时遗憾值目标（式\eqref{eq:regret}）直接映射至单次轨迹推演层面：越早穿过高浓度斑块，函数$m_i(\tau)$越早跃升至高位，积分均值$V_i$也就越高。在给定推演预算下，规划器采用跨周期热启动的交叉熵搜索算法\cite{deboer2005cem}迭代更新航路点$(\mathbf{w}_1, \mathbf{w}_2)$的高斯采样分布，不断向高价值轨迹区域聚焦。

然而，若让交叉熵规划器在全部48个假想海洋上评估并搜索最优路线，同时又用这48个假想海洋上的平均得分来判断“增加算力是否提升了路线质量”，系统将不可避免地陷入严重的“优化器诅咒”\cite{smith2006curse, vanhasselt2010double}。其数学机理在于：有限个假想海洋样本只是真实先验分布的有限蒙特卡洛抽样；当规划推演预算成百上千倍增加时，优化器拥有极高的自由度去精细微调折线转折点，使其恰好擦过这组特定样本中偶然重合的随机正涨落峰值。此时，规划器在内部样本集上观测到的目标函数值会随算力增加而持续虚高，但在未参与搜索的真实海洋中，这些针对样本噪声的微调根本无法带来任何实际收益。

为彻底切断搜索过程与价值评估之间的正向选择偏差，本文借鉴双Q学习中解耦动作选择与价值评估的思想\cite{vanhasselt2010double}，提出搜索—留出双集合解耦架构（图\ref{fig:overview}B）：将48个假想海洋固定划分为两个互不相交的子集——前32个样本构成**搜索集合**$\mathcal{I}_{\text{search}}=\{1,\dots,32\}$，后16个样本构成**留出验证集合**$\mathcal{I}_{\text{hold}}=\{33,\dots,48\}$。在现场同化阶段，全部48个样本均按式\eqref{eq:eakf}同步吸收真实荧光计观测，确保两个子集始终代表相同的后验知识状态；但在轨迹规划阶段，交叉熵优化器仅被允许访问32个搜索样本，通过最大化搜索集均值来选定最优航路点$(\mathbf{w}_1^*, \mathbf{w}_2^*)$：
\begin{equation}
  (\mathbf{w}_1^*, \mathbf{w}_2^*) = \arg\max_{(\mathbf{w}_1, \mathbf{w}_2)} \frac{1}{32}\sum_{i=1}^{32} V_i(\mathbf{w}_1, \mathbf{w}_2)
  \label{eq:cem_search}
\end{equation}
而对于选定的路线$(\mathbf{w}_1^*, \mathbf{w}_2^*)$，其真实的预期收益则完全交由从未参与搜索的16个留出样本进行独立无偏评估：
\begin{equation}
  \hat{V}_{\text{hold}}(\mathbf{w}_1^*, \mathbf{w}_2^*) = \frac{1}{16}\sum_{i=33}^{48} V_i(\mathbf{w}_1^*, \mathbf{w}_2^*)
  \label{eq:hold_eval}
\end{equation}
由于16个留出样本与32个搜索样本共享相同的统计分布及已同化的观测约束，但其未观测区域的随机涨落与搜索集相互独立，式\eqref{eq:hold_eval}能够在完全不接触真实海洋的前提下，精确滤除优化器过拟合样本噪声产生的虚高泡沫。

\subsection{行前假想预演元控制器与对比基线设计}

在解决了单次规划的无偏评估问题后，核心决策难题聚焦于元控制层面：面对特定的机载算力速率、电池容量与当前底图不确定性，AUV究竟应如何在72个决策周期内分配“巡航”、“深度思考”与“停航重采样”？如前所述，经典单步元推理\cite{russell1991right, lin2015metareasoning}试图在每个周期在线比较单步计算价值（由规划器内部改进轨迹外推）与单步信息价值（由假想海洋上的预后验方差缩减估计），但这种短视局部近似既受制于内部搜索轨迹的过思考偏差，又无法预见当前开启GPU消耗的$20\text{ W}$额外功率将在20小时后导致电池提前耗尽、进而丧失最后几小时空间搜索能力的全局长程后果。与此同时，直接用历史任务结果训练策略选择器的做法则因不同任务窗口间高达0.08的自然遗憾值涨落掩盖了0.01量级的策略差异而失效。

既然单步短视评估无法捕捉全周期能量耦合，而跨任务历史学习又受困于环境背景噪声，最直接且严密的破局之道便是在**同一任务窗口内部**、利用航行器自身已校准的假想海洋开展全任务周期的闭环沙盘推演。为此，本文提出**行前假想预演元控制器**（Pre-Mission Rehearsal in Imagination）。首先，提炼出覆盖“行动、思考、停航”基本组合空间的五种代表性时间分配策略构成候选策略集$\mathcal{S}=\{s_1, s_2, s_3, s_4, s_5\}$：
\begin{enumerate}
  \item \textbf{全程巡航（$s_1$）：} 所有周期均保持$1.5\text{ m/s}$巡航采样，仅使用后台低功耗处理器规划（整机功耗$45\text{ W}$）；
  \item \textbf{全程深度思考（$s_2$）：} 所有周期均在巡航采样的同时开启GPU深度思考（整机功耗$65\text{ W}$，基准配置下约15.4小时耗尽电池）；
  \item \textbf{三分之一间歇思考（$s_3$）：} 全程保持巡航采样，每3个周期中开启1个周期的GPU深度思考，其余2个周期使用后台规划（平均功耗约$51.7\text{ W}$，续航约19.3小时）；
  \item \textbf{三分之一停航思考（$s_4$）：} 每3个周期中，有1个周期关停推进器随流漂移并开启GPU深度思考（功耗$35\text{ W}$），其余2个周期巡航并使用后台规划（平均功耗约$41.7\text{ W}$，续航达24小时上限）；
  \item \textbf{三分之一停航重采样（$s_5$）：} 每3个周期中，有1个周期关停推进器随流漂移开展原位密集重采样且仅用后台规划（功耗$15\text{ W}$），其余2个周期巡航并使用后台规划（平均功耗$35\text{ W}$，续航达24小时上限）。
\end{enumerate}

在AUV正式下水出发前（$t=0$，位于母船甲板或水面待命阶段），行前假想预演按如下闭环流程自动选定当前任务的最优策略：从16个留出假想海洋$\mathcal{I}_{\text{hold}}$中抽取$M=4$个样本$\{c^{(j)}\}_{j=1}^4$，将其分别视作假想世界中的“合成真实海洋”。对于每一个合成真实海洋$c^{(j)}$和每一种候选策略$s\in\mathcal{S}$，仿真器调用AUV自身的序列集合调整卡尔曼滤波器、交叉熵规划器、真实的机载基础算力速率$r_{\text{base}}$、真实的GPU加速倍率及真实的$1\text{ kWh}$电池功耗模型，在内部完整推演一遍24小时（72个周期）的闭环巡测任务，计算得出策略$s$在合成真实海洋$c^{(j)}$上的假想随时遗憾值$R_{s,j}$。为最大限度压低4个样本间的蒙特卡洛评估方差，五种候选策略在同一个合成真实海洋$c^{(j)}$上推演时，严格共享完全相同的初始先验集合与随机数种子（即共同随机数方差缩减技术），从而彻底剥离环境背景难度的干扰，纯粹度量策略间的相对优劣。最终，AUV选择在4个留出合成真实海洋上平均假想遗憾值最低的策略$s^*$作为本次实海任务的执行策略：
\begin{equation}
  s^* = \arg\min_{s\in\mathcal{S}} \frac{1}{4}\sum_{j=1}^4 R_{s,j}
  \label{eq:rehearsal}
\end{equation}

为全面剖析本文方法的内在机理，实验部分设置四类对比基线：（1）五种单一固定策略（$s_1$至$s_5$）；（2）**特定硬件开发集最优调参策略（Tuned）**，即针对每一类硬件平台，利用2018至2022年120个开发集任务窗口事后统计出平均随时遗憾值最低的固定策略，代表需要大规模历史任务试错才能获得的先验上限；（3）**经典单步元推理（One-Step Metareasoning）**\cite{russell1991right}，在每个20分钟周期在线对比规划器内部搜索改进率与假想海洋预后验采样收益以动态决定机体与算力动作；（4）**自评预演对照（Rehearsal on Searched Truths）**，其预演流程与式\eqref{eq:rehearsal}完全一致，唯一的区别是从32个搜索集样本$\mathcal{I}_{\text{search}}$（而非留出集$\mathcal{I}_{\text{hold}}$）中抽取4个合成真实海洋，用于检验闭环全任务预演对合成真值来源的敏感性。

\section{实验结果与分析}

\subsection{规划预算扩展下的过思考偏差与留出集合的无偏性验证}

实验首先检验第4.2节提出的第一项核心论断：当单步轨迹规划的推演预算大幅扩展时，规划器自身的内部搜索得分是否会因优化器诅咒而背离真实海洋收益，以及16个留出假想海洋能否提供准确的无偏校准。为此，在开发集抽取的80个赤潮任务窗口上（每个窗口运行3组独立随机种子，共240次实验，装订滞后3天的卫星底图，考察任务前8小时的规划表现），将每次决策的交叉熵搜索推演预算从64次跨越四个数量级逐步扩展至200,000次，同步记录所选路线在32个搜索样本上的规划器内部自评得分、在16个留出假想海洋上的独立评估得分以及在真实海洋中的实际路线价值（图\ref{fig:curse}）。

\begin{figure}[htbp]
  \centering
  \includegraphics[width=\linewidth]{fig_curse}
  \caption{轨迹规划推演预算从64次扩展至200,000次时的优化器诅咒与留出假想海洋无偏校准验证（80个开发集窗口$\times 3$个随机种子，3天滞后底图）。（a）规划路线价值随单次决策推演次数的变化曲线：真实海洋实际价值（黑色实线）与16个留出假想海洋评估值（蓝色虚线）在约2,000次推演后同步趋于平坦，而规划器在32个搜索样本上的内部搜索得分（红色实线）持续陡峭攀升；（b）估计值减去真实海洋价值的评估偏差曲线：内部搜索得分在高预算下产生高达$+0.114$的正向虚高偏差，而基于数据校准不确定性的留出集评估偏差在全预算跨度内始终保持在$0.005$以内；灰色点线展示先验不确定性幅值人为减半（$0.5\times$）或放大（$1.5\times$）时的留出集偏差平移特性。}
  \label{fig:curse}
\end{figure}

如图\ref{fig:curse}a所示，真实海洋中的路线价值随推演预算的增长呈现出显著的边际效益递减与饱和规律：当单次决策推演预算从64次增至2,000次时，真实路线价值从0.749显著提升至0.793（净增$+0.044$），说明在低算力区间增加轨迹搜索确实能够有效规避局部次优、找出穿透高浓度区的走廊；然而，当推演预算从2,000次再激增一百倍至200,000次时，真实海洋路线价值仅从0.793微弱增长至0.803（净增仅$+0.010$），绝大部分真实规划红利在2,000次推演处即已宣告枯竭。与之形成鲜明反差的是，规划器在32个搜索样本上的内部得分（图\ref{fig:curse}a红线）完全无视这一物理饱和边界，从64次推演时的0.762一路加速攀升至2,000次时的0.833，并在200,000次推演时达到惊人的0.917，相对于真实海洋价值产生了高达$+0.114$的虚假乐观偏差（图\ref{fig:curse}b）。这一现象直观印证了连续轨迹规划中的优化器诅咒：当搜索算力过剩时，交叉熵算法通过精巧扭曲折线航路点，专门迎合了32个搜索样本在未观测区域中偶然出现的局部随机正噪点；依赖内部得分的元推理器正是被这一不断膨胀的虚假增量所欺骗，误以为继续投入GPU算力仍在产生丰厚回报。

与之相反，将完全相同的规划路线置于从未参与搜索的16个留出假想海洋上进行评估（图\ref{fig:curse}a蓝线），其得分在64次、2,000次与200,000次推演下分别为0.745、0.794与0.802，与真实海洋价值曲线（0.749、0.793、0.803）近乎完美重合；在横跨四个数量级的全部预算范围内，留出集评估均值与真实海洋实际价值之差始终控制在$0.005$以内（图\ref{fig:curse}b）。更进一步地，当故意破坏先验不确定性校准精度，将生成假想海洋的先验标准差$\sigma_k$人为减半（$0.5\times$）或放大至1.5倍（$1.5\times$）时（图\ref{fig:curse}b灰色点线），留出集评估的绝对水平虽分别产生了$-0.06$与$+0.07$的常数基线偏移，但其对“预算从64次增至200,000次所获相对增益”的估计依然高度稳健：减半与放大不确定性下的留出集评估增益分别为$+0.059$与$+0.054$，与对应工况下的真实海洋实际增益（$+0.054$与$+0.047$）高度吻合。这证明搜索—留出双集合解耦不仅在先验准确校准时具备零偏特性，即便在先验幅值存在一定失配时，依然能够准确捕捉增加算力的边际收益拐点。

\subsection{算力与底图误差耦合相图及停航重采样失效机理分析}

在明确了单步规划的算力饱和规律后，实验进一步考察全任务周期内不同固定分配策略在机载算力速率与底图失真程度二维参数空间中的胜出规律（图\ref{fig:phase}）。通过在后台低功耗处理器速率$r_{\text{base}}\in[5, 200]\text{ 推演/min}$与海真相对于卫星底图偏离标准差$s_{\text{sat}}\in[0.05, 0.30]$构成的网格上开展系统扫描，相图清晰地揭示出具身移动采样中硬件能力与环境不确定性的交互边界\TBD{相图最终扫描完成后的具体分界数值补充}。

\begin{figure}[htbp]
  \centering
  \includegraphics[width=\linewidth]{fig_phase}
  \caption{机载低功耗处理器速率（每分钟5至200次推演）与真实海洋相对于卫星底图偏离程度（0.05至0.30）二维参数空间中的最优固定策略相图。网格颜色标示五种固定时间分配策略中表现最优的策略类别，网格内数值表示该最优策略相对于“全程巡航（Always Go）”策略降低的随时遗憾值收益。}
  \label{fig:phase}
\end{figure}

从图\ref{fig:phase}及测试集定量结果可以提炼出三条清晰的物理规律。第一，**开启高功耗GPU深度思考仅在机载基础算力相对于规划问题过慢时才具有显著正收益**。当AUV搭载慢速后台处理器（$r_{\text{base}}=5\text{ 推演/min}$）时，20分钟周期内的后台常规推演仅能完成100次路线评估，远低于图\ref{fig:curse}a所示的2,000次饱和阈值，导致“全程巡航”策略规划出的路线质量低劣，在2023至2026年测试集上的平均随时遗憾值高达0.374；此时，若每3个周期唤醒1次50倍速的GPU（单周期推演跃升至5,000次，直接跨越饱和拐点），尽管多消耗了部分电能，“三分之一间歇思考”与“全程深度思考”仍将测试集随时遗憾值大幅压低至0.270与0.273，较全程巡航产生高达0.104的显著净收益！反之，当机载后台处理器升级至中速水平（$r_{\text{base}}=20\text{ 推演/min}$）时，单周期后台推演已达400次，配合跨周期热启动机制，数个周期累积的有效搜索深度已逼近饱和区；此时再开启$+20\text{ W}$的GPU不仅无法带来实质性的路线质量跃升，反而因白白消耗电池能量而缩短了总巡航里程，致使“全程巡航”以0.266的随时遗憾值反超“三分之一间歇思考”（0.269）与“全程深度思考”（0.269），成为中速平台上的最优固定策略。

第二，**在巡航途中关停推进器开展停航漂移重采样或停航思考，在所有测试算力与底图误差组合下均得不偿失**。无论是在慢速平台上（“三分之一停航思考”遗憾值0.331，“三分之一停航重采样”高达0.408），还是在中速平台上（两者遗憾值分别为0.328与0.335），包含停航动作的策略均大幅落后于持续巡航策略。这一现象可通过空间信息论与采样信噪比得到严格解释：在空间相关尺度为$\ell=10\text{ km}$（5个2千米网格）、先验合成不确定性$\sigma_3\approx 0.258$的东海赤潮场中，AUV保持$1.5\text{ m/s}$巡航时每20分钟跨越$18\text{ km}$，沿途自动采集18个独立荧光计样本（平均每个2千米网格内采集2个样本，单周期内相关半径内累积9至18个样本）；根据式\eqref{eq:eakf}的卡尔曼方差缩减关系，仅凭巡航途中的走航采样即可将沿途局部后验标准差从0.258迅速压低至$0.1/\sqrt{9}\approx 0.033$以下，局部测量噪声早已不是决策瓶颈。若此时关停推进器驻留原地重采样20分钟，虽能额外获得5个原位样本将同一点位的误差再压低几个毫对数单位，却直接付出了放弃$18\text{ km}$未知海域空间探索的惨重代价，在80千米尺度的任务窗口内严重迟滞了对真正赤潮核心区的搜寻进程。

\begin{figure}[htbp]
  \centering
  \includegraphics[width=\linewidth]{fig_example}
  \caption{典型测试集赤潮巡测任务（2023年4月15日，慢速算力平台$5\text{ 推演/min}$）中“全程巡航”与“三分之一间歇思考”的轨迹与随时遗憾值对比。（a）AUV在下水前装订的3天滞后卫星叶绿素底图；（b）目标日真实海洋叶绿素场及两种策略的24小时航行轨迹，其中蓝色曲线为“全程巡航”轨迹，橙色曲线为“三分之一间歇思考”轨迹（加粗线段标示开启GPU深度思考的20分钟周期），五角星标示真实最高叶绿素浓度水体位置；（c）24小时窗口内AUV向岸基中心实时回报答案的瞬时遗憾值演化曲线，“全程巡航”长期停滞在0.40附近（24小时均值0.43），而“三分之一间歇思考”在约19小时处成功命中全局极值（瞬时遗憾值降至0，24小时均值0.31）。}
  \label{fig:example}
\end{figure}

第三，为直观展示高功耗深度思考在慢速算力平台上扭转战局的物理过程，图\ref{fig:example}剖析了按日期顺序排列的第一个策略差异超过0.1的测试集典型任务实例（2023年4月15日，慢速算力平台$r_{\text{base}}=5\text{ 推演/min}$）。对比图\ref{fig:example}a与图\ref{fig:example}b可见，3天前的滞后卫星底图在窗口东北部与中南部均显示有较高浓度斑块，但经过3天演化与小尺度修正后，真实海洋的全局最高浓度核心（五角星）实际聚集在窗口西南部。在“全程巡航”策略下（蓝色轨迹），由于每周期仅有100次推演预算，交叉熵规划器无力在连续四维航路点空间中探明横跨整个窗口的远距离高价值走廊，导致AUV长期在出发点附近的东北角次优水体中往复盘旋，其上报给岸基中心的最佳浓度估计终生停留在距离真实极值0.40个对数单位之外，24小时时间平均遗憾值高达0.43（图\ref{fig:example}c）。相反，在“三分之一间歇思考”策略下（橙色轨迹），周期性唤醒的GPU（加粗橙色线段）提供了单周期5,000次推演的充裕视野，使规划器在任务前中期成功识别出东北部先验斑块的观测证伪信号，果断规划出一条斜穿任务窗口直插西南部核心区的长距离机动航线，最终在约19小时处精确命中真实最密集赤潮水体，将瞬时遗憾值降至0，全窗口时间平均遗憾值大幅降低至0.31。

\subsection{四类硬件平台上的主实验对比与行前预演有效性分析}

最后，在2023至2026年全部148个测试集赤潮任务窗口上，对五种固定策略、特定硬件开发集最优调参策略（Tuned）、经典单步元推理以及本文提出的行前假想预演控制器（含留出集真值预演与搜索集真值自评预演）开展系统对比。实验涵盖四类典型机载计算硬件平台：（1）慢速处理器平台（$r_{\text{base}}=5\text{ 推演/min}$，GPU额外功耗$+20\text{ W}$）；（2）中速处理器平台（$r_{\text{base}}=20\text{ 推演/min}$，GPU额外功耗$+20\text{ W}$）；（3）快速处理器平台（$r_{\text{base}}=100\text{ 推演/min}$，GPU额外功耗$+20\text{ W}$）；（4）配备高效能专用加速器的中速平台（$r_{\text{base}}=20\text{ 推演/min}$，加速器额外功耗仅$+5\text{ W}$）。全部测试结果汇总于表\ref{tab:main}与图\ref{fig:main}（置信区间通过对赤潮暴发周进行分层自助重采样计算得出）。

\begin{table}[htbp]
  \centering
  \caption{2023至2026年测试集（148个赤潮任务窗口，3天滞后底图，海真偏离0.15，传感器噪声0.1）上各控制策略在四类机载硬件平台上的24小时时间平均随时遗憾值对比（数值越低越好）。括号内数值表示相对于该硬件平台在开发集（120个窗口）上事后选出的最优固定调参策略（标有“调参最优”）的遗憾值差值及基于赤潮周自助法的95\%置信区间。}
  \label{tab:main}
  \small
  \begin{tabular}{lcccc}
    \toprule
    \textbf{控制策略} & \textbf{慢速算力平台} & \textbf{中速算力平台} & \textbf{快速算力平台} & \textbf{高效加速器平台} \\
     & ($5\text{/min}, +20\text{ W}$) & ($20\text{/min}, +20\text{ W}$) & ($100\text{/min}, +20\text{ W}$) & ($20\text{/min}, +5\text{ W}$) \\
    \midrule
    全程巡航（$s_1$） & $0.374\ (+0.104)$ & $\mathbf{0.266}\ (\text{调参最优})$ & \TBD{待填} & \TBD{待填} \\
    全程深度思考（$s_2$） & $0.273\ (+0.003)$ & $0.269\ (+0.003)$ & \TBD{待填} & \TBD{待填} \\
    三分之一间歇思考（$s_3$） & $\mathbf{0.270}\ (\text{调参最优})$ & $0.269\ (+0.003)$ & \TBD{待填} & \TBD{待填} \\
    三分之一停航思考（$s_4$） & $0.331\ (+0.061)$ & $0.328\ (+0.062)$ & \TBD{待填} & \TBD{待填} \\
    三分之一停航重采样（$s_5$） & $0.408\ (+0.138)$ & $0.335\ (+0.069)$ & \TBD{待填} & \TBD{待填} \\
    \midrule
    经典单步元推理\cite{russell1991right} & $0.287\ (+0.017)$ & $0.287\ (+0.021)$ & \TBD{待填} & \TBD{待填} \\
     & {\footnotesize $[+0.009, +0.025]$} & & & \\
    \textbf{行前预演（留出集真值，本文）} & $\mathbf{0.276}\ (+0.007)$ & $\mathbf{0.265}\ (-0.001)$ & \TBD{待填} & \TBD{待填} \\
     & {\footnotesize $[-0.001, +0.015]$} & & & \\
    行前预演（搜索集真值，对照） & $0.273\ (+0.003)$ & $0.268\ (+0.002)$ & \TBD{待填} & \TBD{待填} \\
    \bottomrule
  \end{tabular}
\end{table}

\begin{figure}[htbp]
  \centering
  \includegraphics[width=\linewidth]{fig_main}
  \caption{2023至2026年测试集（148个赤潮任务窗口）上各控制器在四类硬件平台（慢速算力$5\text{/min}$、中速算力$20\text{/min}$、快速算力$100\text{/min}$及低功耗高效加速器$20\text{/min}, +5\text{ W}$）上相对于开发集特定硬件最优调参策略的随时遗憾值差值及95\%自助法置信区间（虚线0基准代表特定硬件事后调参上限，越靠左性能越优）。}
  \label{fig:main}
\end{figure}

表\ref{tab:main}与图\ref{fig:main}的实验数据有力地支撑了本文关于自适应元控制的各项核心结论。首先，**没有任何一种单一固定策略能够跨硬件平台保持最优**。若将中速平台上的最优策略“全程巡航”直接移植至慢速平台，其随时遗憾值将剧增0.104（从0.270恶化至0.374）；反之，若将慢速平台上的最优策略“三分之一间歇思考”固化至中速或更高速平台，则会因不必要的GPU能耗损失续航而落后于全程巡航\TBD{快速算力与高效加速器平台的具体对比数值补充}。这说明传统的固定启发式规则必须针对每一款特定机载处理器和功耗参数组合耗费数年历史任务数据进行事后调参，缺乏跨平台迁移能力。

其次，**经典单步元推理在真实海洋巡测中因严重的过思考偏差而全面落后于调参基线与行前预演**。在慢速与中速算力平台上，经典单步元推理的随时遗憾值均为0.287，分别比对应平台的调参最优策略劣化$+0.017$（95\%置信区间$[+0.009, +0.025]$，统计显著劣化）与$+0.021$。行为轨迹统计揭示了其失效的直接原因：在中速算力平台上，尽管400次后台推演早已足够，经典单步元推理却被规划器在32个搜索样本上持续虚高的内部改进轨迹（图\ref{fig:curse}）所误导，在高达65\%的决策周期内盲目开启GPU进行深度思考，更有28\%的周期选择关停推进电机在原地“停航思考”！这种典型的过度思考病态行为同时浪费了宝贵的推进电能与空间探索时间，证明不经集合解耦与全周期耦合校准的局部单步元推理不适用于受限具身系统。

再次，**本文提出的基于留出集真值的行前假想预演控制器，在无需任何跨任务历史调参的前提下，自动达到了与针对特定硬件事后最优调参策略完全持平的优异性能**。在慢速算力平台上，留出集行前预演自动识别出基础算力的严重不足，高频选出深度思考策略，取得0.276的随时遗憾值，较“全程巡航”大幅降低0.098，与事后调参上限（0.270）之差仅为$+0.007$（95\%置信区间$[-0.001, +0.015]$，统计上无显著差异），且显著优于经典单步元推理（$p < 0.01$）；当硬件切换至中速算力平台时，同一套行前预演算法无需修改任何一行参数，自动感知到后台算力已趋充裕而GPU额外功耗会导致提前断电，从而果断转向以巡航为主的策略，取得0.265的随时遗憾值（甚至较开发集调参最优策略0.266微弱领先$-0.001$，较单步元推理大幅领先0.022）。在全部测试任务中，行前预演几乎从不选择停航类策略（选中比例仅为1\%），准确复现了真实海洋中的策略相图规律\TBD{快速算力与高效加速器平台上行前预演的具体数值表现补充}。

最后，对比表\ref{tab:main}最后两行可以发现一个极具启发性的现象：**在闭环全任务行前预演中，即便从搜索集（而非留出集）中抽取合成真实海洋（自评预演对照），依然取得了与留出集预演高度一致的优异表现**（慢速平台0.273，中速平台0.268）。为何单步评估时搜索集会产生高达$+0.114$的灾难性过拟合（图\ref{fig:curse}），而在24小时闭环行前预演中却能保持准确的策略排序？深入剖析预演内部数据发现，在全任务闭环预演中，一旦某个假想海洋$c^{(j)}$被抽作“合成真实海洋”，它在整场24小时模拟任务中便扮演固定的外部物理环境角色，而模拟任务内部的规划器在每一步是通过同化带噪观测来搜索路径，无法直接针对未知的未来观测实现作弊；同时，由于假想海洋由高斯随机场生成，其极值尖峰分布与真实海洋存在微弱的高阶矩差异，导致预演计算出的绝对假想遗憾值在所有策略上均比真实海洋遗憾值系统性偏高0.05至0.07个对数单位。然而，由于五种策略在预演中严格共享相同的合成真实海洋与随机数序列（共同随机数技术），这一0.05至0.07的系统性加性常数偏移在五种策略上同向发生，并在按式\eqref{eq:rehearsal}做策略间作差比较时完美相互抵消！例如在中速算力平台上，留出集预演与搜索集预演估计的“思考策略减去巡航策略”相对遗憾值差值分别为$-0.002$与$+0.006$，均与真实海洋中的实际策略差值（$+0.002$）紧密吻合。这表明，全周期闭环预演结合共同随机数差分比较，天然具备对绝对评估基线偏移的免疫能力，进一步夯实了行前假想预演在实际工程部署中的鲁棒性。

\section{结论与展望}

本文围绕携带海洋大数据先验底图的自主水下航行器在动态巡测中如何权衡“深入思考”、“巡航采样”与“停航重采样”这一核心科学问题，依托2018至2026年东海3197天多源融合卫星叶绿素大数据与422个春季赤潮应急巡测任务窗口，开展了从大数据不确定性量化、集合同化建模、规划器过思考机理剖析到自适应元控制器设计的全链条研究。研究得出三项主要结论：（1）在基于有限假想海洋集合的连续轨迹规划中，随着推演算力跨越数个数量级增长，规划器内部搜索得分在真实收益早已饱和后仍因优化器诅咒而持续攀升，而将假想海洋解耦为搜索集与留出验证集可将全预算范围内的评估偏差严格控制在0.005以内；（2）深度时空神经网络无法仅凭日均平滑遥感底图确定性外推小尺度赤潮斑块的日际演变，高功耗GPU深度思考仅在机载基础算力不足以覆盖单周期有效规划深度时产生显著净收益，而在相关尺度为十千米级的中尺度场中，巡航途中关停推进器开展停航重采样始终得不偿失；（3）在出发前利用假想海洋开展24小时全周期闭环预演并结合共同随机数差分择优，能够有效克服经典单步元推理的过思考病态与离线学习的低信噪比困境，在无需特定硬件历史调参的条件下自动匹配各算力平台的事后最优调参性能。

本研究在海流运动表征与多平台协同方面仍存在一定局限：当前模型在随流漂移参考系下假设表层水体在80千米窗口内作准刚性平流，尚未显式刻画强锋面区的水平剪切变形及藻华昼夜垂直迁移引起的三维垂向分布差异，且行前预演在任务中途遭遇极端突发水文突变时缺乏在线重预演修正机制。未来的研究工作将沿着两个方向拓展：一是将三维斜压海流数值预报集合与卫星水色大数据融合，在时变剪切流场中研究AUV三维锯齿剖面机动与机载实时科学计算的联合资源分配；二是将单机行前假想预演拓展至空海潜跨域异构无人集群（如无人艇与多台AUV协同），探索通信受限条件下分布式想象推演与信息共享的最优权衡准则。
